\clearpage
\pdfbookmark[1]{Nachrichten. 30. Juli 1849.}{anzeige_1849_122_124}
\SourceLeaf{122}
\par\noindent\hfill\rule{0.30\linewidth}{0.4pt}\hfill\mbox{}\par
\begin{center}
\textbf{Nachrichten von der Königlichen Gesellschaft der Wissenschaften zu Göttingen. 30. Juli 1849.}
\end{center}

\par\noindent\hfill\rule{0.30\linewidth}{0.4pt}\hfill\mbox{}\par

Die Vorlesung, die der Geheime Hofrat \textsc{Gauss} in der Sitzung der Königlichen Gesellschaft der Wissenschaften am 16. Juli hielt, betraf die
\begin{center}
\textit{Theorie der algebraischen Gleichungen.}
\end{center}
Von der ausführlichen Abhandlung, die in Kürze im Druck erscheinen wird, geben wir hier nur den folgenden Bericht.

Sie besteht aus zwei Abteilungen, von denen jede einem besonderen Gegenstand gewidmet ist.

Die erste Abteilung behandelt den Fundamentalsatz der Theorie der algebraischen Gleichungen, nämlich die Zerlegbarkeit jedes Polynoms in einer Variablen in Faktoren. Darüber hatte der Verfasser bereits 1799 eine Abhandlung unter dem Titel \textit{Ein neuer Beweis des Satzes, dass jedes Polynom} \textit{in einer Variablen in reelle Faktoren ersten oder zweiten Grades zerlegt werden kann} veröffentlicht. Diese Abhandlung hatte einen doppelten Zweck: erstens zu zeigen, dass sämtliche bis dahin versuchten Beweise des ausgesprochenen Satzes unzureichend und trügerisch waren, und zweitens einen vollkommen strengen neuen Beweis zu geben. Auf den ersten Gegenstand noch einmal zurückzukommen, wäre unnötig. Es muss jedoch bemerkt werden, dass der Verfasser jenem ersten neuen Beweis später noch zwei weitere folgen ließ, die sich im 3. Band
\SourceLeaf{123}
der \textit{Commentationes recentiores} unserer Gesellschaft befinden, und dass ein vierter zuerst von \textsc{Cauchy} aufgestellt wurde. Alle diese vier vollkommen strengen Beweise beruhen auf ebenso vielen unterschiedlichen Grundlagen. Gemeinsam ist ihnen jedoch, dass sie alle \textit{unmittelbar} nur die Existenz eines Faktors der betreffenden Funktion nachweisen. Die Strenge der Beweise wird dadurch nicht beeinträchtigt. Denn nach Abtrennung dieses einen Faktors von der gegebenen Funktion bleibt offensichtlich eine ähnliche Funktion niedrigeren Grades übrig, auf die der Satz erneut angewandt werden kann. Die fortgesetzte Wiederholung des Verfahrens ergibt schließlich eine vollständige Zerlegung der ursprünglichen Funktion in Faktoren der bezeichneten Art. Jede Beweisführung gewinnt jedoch offenbar eine höhere Vollendung, wenn gezeigt wird, dass sie die Existenz sämtlicher einzelner Faktoren auf einmal unmittelbar anschaulich machen kann. Dies ist beim ersten Beweis tatsächlich der Fall, wie in der erwähnten Abhandlung (Art. 23) bereits angedeutet, aber nicht weiter ausgeführt wurde. Die vorliegende Abhandlung gibt nun diese Erweiterung, jedoch als Bestandteil einer Umarbeitung der Abhandlung von 1799 und nicht als deren Zusatz. Der Beweis wird hier selbstständig neu dargestellt, aber mit einigen wesentlichen Änderungen im Ablauf, die eine bedeutende Vereinfachung ergeben. Auch die gesamte Einkleidung ist verändert. Sie war in der älteren Schrift bewusst so gewählt, dass jede Einbeziehung imaginärer Größen vollständig vermieden werden konnte. Die Beweisführung steht jedoch sowohl ihrem \textit{Ursprung} als auch ihrem wahren Wesen nach in engstem Zusammenhang mit der Vorstellung komplexer Größen. Seit diese sich in der Mathematik eingebürgert hat, ist dies aufmerksamen Lesern jener Abhandlung nicht entgangen. Der Verfasser hielt es daher für angemessen, jene erste Einkleidung jetzt aufzugeben und die Argumentation auf ihr eigentliches Feld zurückzuführen, dessen Grundbegriffe nun jedem geläufig sind, zumal dadurch nicht nur größere Einfachheit, sondern auch noch etwas größere Allgemeinheit gewonnen wird.

Die zweite Abteilung ist den algebraischen Gleichungen mit drei Gliedern gewidmet. Ihre Besonderheit ist, dass einige Methoden zur numerischen Lösung von Gleichungen bei ihnen eine Geschmeidigkeit und Eleganz erreichen können, von der ihre Anwendung auf Gleichungen weniger einfacher Gestalt weit entfernt bleibt. Dies gilt insbesondere für die Lösung
\SourceLeaf{124}
durch unendliche Reihen und für die indirekte Methode. Die Entwicklung dieser Methoden für Gleichungen jener Form scheint daher umso mehr eine besondere Darstellung zu verdienen, als solche Gleichungen tatsächlich sehr häufig vorkommen.

Die Lösung durch unendliche Reihen hat der Verfasser jedoch aus seinem gegenwärtigen Vorhaben vollständig ausgeschlossen. Er hat lediglich bemerkt, dass sich für \textit{jede} Wurzel einer solchen Gleichung, ob reell oder imaginär, eine konvergente Reihe finden lässt, die nach einem leicht erkennbaren Gesetz fortschreitet. So schön diese Lösungsart in allgemein theoretischer Hinsicht auch ist, wird bei tatsächlichen praktischen Anwendungen doch in allen Fällen den indirekten Methoden der Vorzug gegeben, in denen jene Konvergenz nicht sehr schnell ist.

Diese indirekten Methoden sind nun der Gegenstand der zweiten Abteilung. Es handelt sich hier um \textit{zwei} Methoden, denn das Verfahren zur Bestimmung der imaginären Wurzeln unterscheidet sich völlig von dem für die reellen Wurzeln. Vom letzteren hat der Verfasser bereits bei anderen Gelegenheiten einige Beispiele für besondere Fälle gegeben und dabei zugleich die allgemeine Anwendbarkeit des Verfahrens angedeutet. Obwohl diese Verallgemeinerung keinerlei Schwierigkeit bereitet, glaubte der Verfasser doch, dass eine Zusammenstellung gebrauchsfertiger Vorschriften willkommen wäre, zumal die Zahl der dabei zu unterscheidenden und einzeln zu behandelnden Fälle beträchtlich ist.

Die Bestimmung der imaginären Wurzeln auf indirektem Weg ist viel schwieriger als die der reellen, sofern nicht schon einigermaßen angenäherte Werte von anderer Seite bekannt sind. Denn jene müssen in einem unendlichen Gebiet von zwei Dimensionen gesucht werden, diese nur in einem unendlichen Gebiet von einer Dimension. Bei Gleichungen mit drei Gliedern lässt sich diese Schwierigkeit durch einen einfachen, sehr naheliegenden Kunstgriff umgehen, der jedoch bisher, wie es scheint, sonst noch nicht benutzt wurde. Auch für diese Aufgabe sind die zur Lösung erforderlichen Vorschriften vollständig und in gebrauchsfertiger Form mitgeteilt.
\par\noindent\hfill\rule{0.30\linewidth}{0.4pt}\hfill\mbox{}\par
